\section*{Conclusiones}

El estudio permitió constatar que la arquitectura por capas sigue siendo un patrón fundamental dentro de la ingeniería de software, no por su rigidez, sino por su capacidad de adaptación a contextos tecnológicos diversos. Su persistencia en el tiempo demuestra que los principios de modularidad, separación de responsabilidades y jerarquía funcional continúan siendo esenciales para el desarrollo de sistemas sostenibles.

En primer lugar, los resultados evidencian que el patrón por capas ha evolucionado hacia modelos híbridos, en los que convive con paradigmas como la Arquitectura Limpia, los microservicios y los pipelines autocontenidos. Esta integración refuerza la autonomía de los módulos y mejora la escalabilidad sin perder la coherencia conceptual del diseño.

En segundo lugar, la literatura analizada revela que la mantenibilidad y la sostenibilidad arquitectónica son los desafíos más críticos. La erosión estructural, causada por violaciones de dependencias y falta de disciplina organizacional, puede neutralizar las ventajas del patrón. En consecuencia, la adopción de herramientas de análisis estático, visualización y refactorización automatizada se consolida como una práctica necesaria para prolongar la vida útil de los sistemas.

Asimismo, se comprobó que el rendimiento no debe considerarse una limitación inherente a las capas, sino una variable controlable mediante una adecuada distribución de responsabilidades y la aplicación de modelos de evaluación temprana. La incorporación de técnicas de modelado, como las redes de colas o los algoritmos de clustering, permite anticipar cuellos de botella y optimizar la estructura antes de su implementación.

Finalmente, el impacto práctico del patrón por capas trasciende la ingeniería técnica: influye directamente en la organización del trabajo, la comunicación entre equipos y la capacidad de evolución del software. Su futuro dependerá de la convergencia entre diseño arquitectónico, automatización inteligente y sostenibilidad estructural. En este sentido, la arquitectura por capas continúa siendo no solo una referencia histórica, sino una herramienta vigente para enfrentar los retos de la complejidad creciente en el desarrollo de software.
